\documentclass[12pt,oneside]{book}

% ==========================================================
% METADATOS
% ==========================================================
\newcommand{\universidad}{Universidad de Buenos Aires (UBA)}
\newcommand{\materia}{Derechos Reales}
\newcommand{\titulo}{Conjuntos Inmobiliarios}
\newcommand{\autor}{Isaac Edgar Camacho Ocampo}
\newcommand{\anio}{2025}

% ==========================================================
% PAQUETES
% ==========================================================
\usepackage[utf8]{inputenc}
\usepackage[T1]{fontenc}
\usepackage[spanish,es-nodecimaldot]{babel}

\usepackage[
    top=2.5cm,
    bottom=2.5cm,
    left=3cm,
    right=2.5cm,
    headheight=15pt
]{geometry}

\usepackage{amsmath}
\usepackage{amssymb}
\usepackage{mathtools}

\usepackage{graphicx}
\usepackage{float}
\usepackage{caption}
\usepackage{subcaption}

\usepackage{booktabs}
\usepackage{array}
\usepackage{longtable}
\usepackage{tabularx}

\usepackage{enumitem}

\usepackage{xcolor}
\usepackage[most]{tcolorbox}

\usepackage{fancyhdr}
\usepackage{ragged2e}

\usepackage[
    colorlinks=true,
    linkcolor=blue!50!black,
    citecolor=green!40!black,
    urlcolor=blue!60!black,
    pdftitle={\titulo},
    pdfauthor={\autor},
    pdfsubject={Apuntes universitarios},
    bookmarks=true
]{hyperref}

% ==========================================================
% CONFIGURACIÓN
% ==========================================================
\graphicspath{
    {./img/}
    {./graficos/}
    {./figuras/}
}

\setcounter{tocdepth}{2}

\setlength{\parindent}{0pt}
\setlength{\parskip}{0.5em}

\setlist[itemize]{
    topsep=0.3em,
    itemsep=0.2em
}

\setlist[enumerate]{
    topsep=0.3em,
    itemsep=0.2em
}

\clubpenalty=10000
\widowpenalty=10000

% ==========================================================
% COMANDOS
% ==========================================================
\newcommand{\abs}[1]{\left|#1\right|}
\newcommand{\norm}[1]{\left\lVert#1\right\rVert}
\newcommand{\vect}[1]{\mathbf{#1}}
\newcommand{\deriv}[2]{\frac{d#1}{d#2}}
\newcommand{\pderiv}[2]{\frac{\partial#1}{\partial#2}}

% ==========================================================
% ENTORNOS / CAJAS
% ==========================================================
\newtcolorbox{definition}[1]{
    enhanced,
    breakable,
    title={Definición: #1},
    fonttitle=\bfseries,
    colback=blue!3,
    colframe=blue!50!black,
    boxrule=0.8pt,
    arc=2mm
}

\newtcolorbox{important}{
    enhanced,
    breakable,
    title={Importante},
    fonttitle=\bfseries,
    colback=yellow!8,
    colframe=orange!70!black,
    boxrule=0.8pt,
    arc=2mm
}

\newtcolorbox{example}[1]{
    enhanced,
    breakable,
    title={Ejemplo: #1},
    fonttitle=\bfseries,
    colback=green!4,
    colframe=green!40!black,
    boxrule=0.8pt,
    arc=2mm
}

\newtcolorbox{formula}[1]{
    enhanced,
    breakable,
    title={Fórmula / Regla: #1},
    fonttitle=\bfseries,
    colback=purple!4,
    colframe=purple!50!black,
    boxrule=0.8pt,
    arc=2mm
}

\newtcolorbox{exercise}[1]{
    enhanced,
    breakable,
    title={Ejercicio: #1},
    fonttitle=\bfseries,
    colback=gray!5,
    colframe=gray!60!black,
    boxrule=0.8pt,
    arc=2mm
}

\newtcolorbox{note}{
    enhanced,
    breakable,
    title={Nota},
    fonttitle=\bfseries,
    colback=white,
    colframe=gray!50,
    boxrule=0.6pt,
    arc=2mm
}

% ==========================================================
% CONFIGURACIÓN FINAL
% ==========================================================
\pagestyle{fancy}
\fancyhf{}
\fancyhead[L]{\nouppercase{\leftmark}}
\fancyhead[R]{\materia}
\fancyfoot[C]{\thepage}
\renewcommand{\headrulewidth}{0.4pt}
\renewcommand{\footrulewidth}{0pt}

\fancypagestyle{plain}{
    \fancyhf{}
    \fancyfoot[C]{\thepage}
    \renewcommand{\headrulewidth}{0pt}
}

% ==========================================================
\begin{document}

% ==========================================================
% PORTADA
% ==========================================================
\begin{titlepage}

\thispagestyle{empty}

\begin{center}

\vspace*{1cm}

{\large \textsc{\universidad}}

\vspace{3.5cm}

{\Huge \textbf{\materia}}

\vspace{1cm}

{\LARGE \titulo}

\vspace{0.8cm}

\textit{Estudiar con constancia es el camino para comprender en profundidad.}

\vspace{1.2cm}

\rule{0.70\textwidth}{0.5pt}

\vspace{0.5cm}

{\large Apuntes de estudio}

\vspace{0.5cm}

\rule{0.70\textwidth}{0.5pt}

\vfill

{\large \autor}

\vspace{0.3cm}

{\large \anio}

\end{center}

\vfill

\begin{flushleft}
\textit{Este es un modesto aporte a los estudiantes de ``Derechos Reales'' no pretende sustituir el material oficial de la materia.}
\end{flushleft}

\end{titlepage}

% ==========================================================
% ÍNDICE
% ==========================================================
\tableofcontents

% ==========================================================
% CONTENIDO
% ==========================================================

\chapter*{Presentación}
\addcontentsline{toc}{chapter}{Presentación}
\markboth{Presentación}{Presentación}

Estos apuntes abordan el surgimiento, la problemática jurídica y la regulación definitiva de los \textbf{conjuntos inmobiliarios} en el derecho argentino. Se analiza la evolución histórica desde su organización bajo \textbf{derechos personales} ---con sus graves riesgos para los adquirentes--- hasta su incorporación como \textbf{derecho real de propiedad horizontal especial} en el Código Civil y Comercial de la Nación (Ley 26.994). Se estudian los artículos 2073 a 2086, la distinción entre partes comunes y privativas, el reglamento, las facultades y restricciones de los propietarios y el complejo proceso de adecuación de los emprendimientos preexistentes.
% TODO: contenido incompleto o ambiguo en las notas originales (las notas anuncian los arts. 2073 a 2086, pero en el desarrollo sólo se transcriben los arts. 1887, 2073, 2074, 2075 tercer párrafo, 2077 y 2079)

\section*{Utilidad para la vida profesional}

Quien ejerza como abogado, escribano o asesor inmobiliario se enfrentará con frecuencia a operaciones sobre inmuebles en countries, barrios cerrados, parques industriales o clubes náuticos. Saber bajo qué régimen está organizado cada emprendimiento ---derechos reales o derechos personales--- determina qué derechos tiene el cliente, qué riesgos enfrenta, si puede hipotecar el bien, afectarlo a bien de familia, o qué ocurre si el desarrollista quiebra. Este conocimiento es, además, indispensable para redactar o revisar reglamentos, escrituras traslativas de dominio y contratos relacionados con este tipo de propiedad. Proteger al cliente en el ámbito de la propiedad inmobiliaria contemporánea exige dominar esta materia.

\section*{Hilo conductor de los temas}

Antes del Código Civil y Comercial, quien compraba una casa en un country no adquiría un derecho de propiedad real sobre el terreno: adquiría acciones de una sociedad anónima que era la titular de toda la tierra. Si esa sociedad quebraba, el adquirente podía perderlo todo. Esa inseguridad jurídica es el punto de partida (Capítulo~1). Las provincias intentaron soluciones parciales e inconstitucionales y la doctrina debatió el tema (Capítulo~2). Finalmente, el Código creó un régimen propio: el derecho real de propiedad horizontal especial (Capítulo~3), con reglas claras sobre qué corresponde a cada propietario, qué se comparte y cómo funciona la comunidad (Capítulo~4).

% ----------------------------------------------------------
\chapter[Marco histórico]{Marco histórico y surgimiento de los conjuntos inmobiliarios}

\section{Qué son los conjuntos inmobiliarios}

Bajo la denominación jurídica de \textbf{conjuntos inmobiliarios}, el Código Civil y Comercial ha regulado distintas formas de organización que eran habituales y se habían desarrollado a lo largo de las últimas décadas bajo las denominaciones de barrios cerrados, countries y clubes de campo.

En todos los casos se trata de \textbf{nuevas manifestaciones de propiedad inmobiliaria} que tuvieron un amplio desarrollo en todo el territorio nacional, incluso antes de que existiera una ley nacional que las regulara.

\section{La organización bajo derechos personales}

La mayoría de estos emprendimientos se organizaba en el ámbito de los \textbf{derechos personales}. El mecanismo era el siguiente: la \textbf{tierra madre} ---aquella en la que se desarrolla y se organiza el complejo--- pertenecía en dominio a una persona jurídica, que tenía el dominio absoluto, exclusivo, perpetuo y excluyente.

Lo que adquiría el comprador de un inmueble en el emprendimiento \textbf{no era un derecho real, sino una acción}. Como en materia de derechos personales prima la \textbf{autonomía de la voluntad}, se regulaba que esa acción le confiriera derechos sobre un determinado lote.

Muchos emprendimientos se desarrollaron sin dificultad bajo este esquema, pero otros generaron graves problemas.

\section{Los riesgos: subastas y quiebras}

El sistema de derechos personales reveló su fragilidad en distintos casos:

\begin{itemize}
    \item Casos en los que, a través de una \textbf{ejecución individual}, se subastó la tierra madre, es decir, la que pertenecía en propiedad a la persona jurídica (sociedad anónima).
    \item Casos en los que la persona jurídica \textbf{quebró por deudas}.
\end{itemize}

Cuando se subasta la tierra, se subasta \textbf{con lo plantado y lo construido}.

\begin{example}{Quiebra de la sociedad titular de la tierra}
Si una sociedad anónima era titular de toda la tierra de un country y quebraba, el bien se subastaba íntegramente, incluyendo las construcciones de cada propietario. Los compradores de lotes perdían todo porque no tenían derechos reales sobre la tierra, sino solo acciones de la empresa.
\end{example}

% ----------------------------------------------------------
\chapter[Regulación provincial y debate doctrinal]{Regulación provincial, inconstitucionalidad y debate doctrinal}

\section{El intento de las provincias: la norma bonaerense}

Ante la ausencia de una ley nacional, diversas provincias intentaron regular la situación. La Provincia de Buenos Aires dictó una norma inconstitucional ---o, al menos, susceptible de ser declarada inconstitucional--- que establecía que los nuevos emprendimientos debían organizarse \textbf{afectándolos a propiedad horizontal}.

La inconstitucionalidad tenía dos fundamentos:

\begin{itemize}
    \item Las provincias \textbf{carecen de competencia} para legislar en materia de propiedad, reservada al Congreso Nacional.
    \item La propiedad horizontal clásica implica que el \textbf{terreno es bien común} de todos los propietarios, lo cual es incompatible con la naturaleza de un country.
\end{itemize}

La Provincia de Buenos Aires era consciente de la inconstitucionalidad y de que la propiedad horizontal implica la propiedad común de todo el terreno; sin embargo, no encontraba otro modo de aportar algo de seguridad a este tipo de emprendimientos. Por ello, \textbf{la sanción de una ley nacional era imprescindible}.

\section{Por qué la propiedad horizontal clásica no servía}

Existe una contradicción esencial entre la propiedad horizontal y los countries. Lo que es común en la propiedad horizontal es el \textbf{terreno}. Además, se trata de edificios ya construidos. Nada de esto ocurre en un country, pues quien adquiere un inmueble en un barrio cerrado o en un country difícilmente desee compartir con todos sus vecinos el terreno en el que construirá su inmueble.
% TODO: contenido incompleto o ambiguo en las notas originales (la transcripción dice "edificios en construcción ya construidos")

\begin{note}
La Dra. Marina Mariani de Vidal estudió durante años esta problemática y caracterizó a los inmuebles construidos en countries organizados bajo propiedad horizontal como \textbf{``grandes castillos construidos sobre arena''}, en referencia a la inseguridad jurídica que los rodeaba.
\end{note}

\section{El debate doctrinal}

\subsection{Postura en favor de los derechos reales}

Gran parte de la doctrina consideró que todos estos emprendimientos representan \textbf{nuevas manifestaciones de propiedad inmobiliaria}, de modo que se imponía legislarlos en el ámbito de los \textbf{derechos reales}, ya sea mediante:

\begin{itemize}
    \item la creación de un \textbf{derecho real autónomo}, o bien
    \item un \textbf{derecho real de propiedad horizontal especial}, que comparta todas las cosas comunes (espacios y servicios comunes) y respete esencialmente la propiedad individual del terreno de cada propietario.
\end{itemize}

Cuando se dio a conocer el proyecto de Código Civil y Comercial, este incluía en el artículo 1887 a los conjuntos inmobiliarios como derechos reales, lo cual generó gran satisfacción en la doctrina. Sin embargo, el artículo 2075 del proyecto original admitía la posibilidad de organizarlos también como \textbf{derechos personales o mixtos}. Esto implicaba que quedaba bajo la decisión del emprendedor o desarrollista someterlos a uno u otro régimen, idea que no fue compartida en la postura aquí reseñada.

\subsection{La ponencia presentada en el Quinto Congreso de Derecho Privado}

En la ponencia presentada en el Quinto Congreso de Derecho Privado se rechazó rotundamente la posibilidad de que algunos conjuntos inmobiliarios pudieran estructurarse como derechos personales. Al emprendedor o urbanizador le conviene económicamente desarrollar el emprendimiento bajo derechos personales, porque ese ámbito es \textbf{más flexible y menos oneroso}: no exige cumplir con todos los requisitos propios de los derechos reales, a saber:

\begin{itemize}
    \item plano aprobado;
    \item escritura pública de cada uno de los propietarios;
    \item inscripción registral para la oponibilidad a terceros.
\end{itemize}

\subsection{La perspectiva del adquirente frente a la del desarrollista}

Se distingue entre lo que busca el desarrollista y lo que necesita el adquirente. Apelando a la experiencia profesional en la zona norte del Gran Buenos Aires, el adquirente atiende a: la ubicación, el precio, el tipo de construcción, los servicios disponibles, la distancia a las lagunas (característica de estos emprendimientos), el tiempo de traslado al trabajo y la cercanía de un colegio para los hijos. Son cuestiones que evalúan sus posibilidades económicas, sus deseos y sus aspiraciones, y en ellas \textbf{el derecho debe guardar un respetuoso silencio}, porque pertenecen a la esfera personal.

El adquirente difícilmente pueda comprender las implicancias jurídicas de estar bajo derechos personales o reales; esa función protectora corresponde al derecho. No es tarea del desarrollista cautelar el orden público ni el interés social de la propiedad: esa es la \textbf{función del derecho real}. El derecho regula la propiedad considerando, más allá del interés inmediato del comprador:

\begin{itemize}
    \item la \textbf{finitud humana} (de allí el régimen de sucesiones);
    \item la \textbf{oponibilidad erga omnes};
    \item la extensión del derecho más allá de la vida del adquirente.
\end{itemize}

\subsection{Ventajas exclusivas del derecho real para el propietario}

Solo el titular de un derecho real puede ejercer los siguientes derechos, imposibles bajo el régimen de derechos personales:

\begin{itemize}
    \item afectar el inmueble al \textbf{régimen de protección de la vivienda familiar};
    \item \textbf{derecho de habitación vitalicio y gratuito} del cónyuge o conviviente supérstite;
    \item \textit{ius persequendi} y preferente (derecho de perseguir la cosa y de ser preferido frente a otros acreedores);
    \item posibilidad de constituir nuevos derechos reales, como usufructo e hipoteca;
    \item \textbf{oponibilidad erga omnes} (frente a todos) mediante inscripción registral.
\end{itemize}

% ----------------------------------------------------------
\chapter[Regulación en el CCyC]{Regulación en el Código Civil y Comercial}

\section{Incorporación como derecho real: artículos 1887 y 2075}

El nuevo cuerpo legislativo incorpora los conjuntos inmobiliarios en el artículo 1887 como \textbf{derechos reales}. Durante el tránsito del proyecto a la sanción definitiva se modificó el artículo 2075, eliminándose la posibilidad de constituir conjuntos inmobiliarios como derechos personales.

\begin{quote}
\textbf{Artículo 1887 CCyC (inciso aplicable).} \textit{Los conjuntos inmobiliarios quedan incorporados al catálogo de derechos reales del Código Civil y Comercial de la Nación. A partir de la sanción del Código (31 de agosto de 2015), todo conjunto inmobiliario debe ser organizado exclusivamente como derecho real.}
\end{quote}
% TODO: contenido incompleto o ambiguo en las notas originales (el texto atribuido al art. 1887 figura así en las notas y no se contrastó con la fuente)

\begin{quote}
\textbf{Artículo 2075, tercer párrafo, CCyC.} \textit{Todos los conjuntos inmobiliarios preexistentes a la sanción del nuevo ordenamiento que se hubieran establecido como derechos personales, o en los que coexistan derechos personales y derechos reales, se deben adecuar a las previsiones normativas que regulan este derecho real.}
\end{quote}

\begin{important}
El Código no solo impide la constitución como derechos personales ---todo conjunto inmobiliario a partir del 31 de agosto de 2015 debe ser constituido única y exclusivamente como derecho real---, sino que también \textbf{impone la adecuación de los anteriores}.
\end{important}

\section{Tipos de conjuntos inmobiliarios: artículo 2073}

\begin{quote}
\textbf{Artículo 2073 CCyC.} \textit{Son conjuntos inmobiliarios los clubes de campo, barrios cerrados o privados, parques industriales, empresariales o náuticos, y todo otro emprendimiento urbanístico independientemente del destino de vivienda permanente o temporaria, laboral, comercial o empresarial que tenga, comprendidos asimismo aquellos que contemplan usos mixtos.}
\end{quote}

No es el destino del emprendimiento lo que lo convierte en conjunto inmobiliario, sino los \textbf{elementos característicos} regulados en el artículo 2074: cerramiento perimetral, partes privativas y partes comunes interdependientes, que forman un todo no escindible.

\section{Elementos característicos: artículo 2074}

\begin{quote}
\textbf{Artículo 2074 CCyC.} \textit{Son elementos característicos de estas urbanizaciones, los siguientes, sin que su enumeración sea taxativa: cerramiento, partes comunes y privativas, estado de indivisión forzosa y perpetua de las partes, lugares y bienes comunes, reglamento por el que se establecen órganos de funcionamiento, limitaciones y restricciones a los derechos particulares y régimen disciplinario, obligación de contribuir con los gastos y cargas comunes y entidad con personería jurídica que agrupe a los propietarios de las unidades privativas.}
\end{quote}

\section{La obligación de adecuación y sus dificultades prácticas}

El artículo 2075 ordena la adecuación de los emprendimientos preexistentes, pero \textbf{no establece un plazo} para que los organizados como derechos personales o mixtos se adecuen, \textbf{ni una sanción} para quienes no lo hagan. En consecuencia, durante muchos años coexistirán conjuntos inmobiliarios anteriores a 2015, regulados como derechos personales y que siguen funcionando así, y otros posteriores, regulados como derechos reales.

La adecuación es compleja porque requiere:

\begin{itemize}
    \item planos de mensura aprobados;
    \item reglamento de propiedad;
    \item escritura pública de afectación;
    \item inscripción registral para cada unidad.
\end{itemize}

La conversión a derechos reales es más conveniente y ofrece mayores seguridades, pero el problema radica en los \textbf{elevadísimos costos}, lo que explica que la ley no haya impuesto plazo ni sanción.

% ----------------------------------------------------------
\chapter[Estructura y funcionamiento]{Estructura y funcionamiento del conjunto inmobiliario}

\section{Propiedad horizontal especial}

Los conjuntos inmobiliarios quedan regulados como una \textbf{propiedad horizontal especial}. La diferencia fundamental con la propiedad horizontal clásica es la siguiente: lo esencialmente común en la propiedad horizontal es el terreno; en cambio, conforme al artículo 2077, en el conjunto inmobiliario cada propietario tiene la \textbf{propiedad exclusiva del lote}. Así como en la propiedad horizontal el derecho recae sobre un ``cubo de aire'', aquí el propietario tiene la propiedad exclusiva de la tierra: existen \textbf{lotes} como objeto de esta propiedad.

\begin{quote}
\textbf{Artículo 2077 CCyC (unidad funcional).} \textit{El derecho de propiedad del propietario en un conjunto inmobiliario se ejerce sobre unidades funcionales privativas, cuya superficie se destina a uso privado. La unidad funcional puede encontrarse construida o en proceso de construcción; debe contar con independencia funcional vinculada a su destino y salida a la vía pública por vías o elementos comunes o en forma directa.}
\end{quote}

También difieren los requisitos de afectación. Para afectar un inmueble a propiedad horizontal debe existir un edificio construido y la propiedad común del terreno. Para afectarlo a conjunto inmobiliario, el inmueble puede estar construido o en proceso de construcción, y el terreno es de propiedad exclusiva del titular. Por eso se habla de una \textbf{propiedad horizontal especial}.

\subsection{Cuadro comparativo: propiedad horizontal clásica y conjunto inmobiliario}

\begin{table}[H]
\centering
\small
\renewcommand{\arraystretch}{1.3}
\begin{tabularx}{\textwidth}{>{\raggedright\arraybackslash}X >{\raggedright\arraybackslash}X}
\toprule
\textbf{Propiedad horizontal clásica} & \textbf{Conjunto inmobiliario} \\
\midrule
El terreno es siempre bien común de todos los copropietarios. & Cada propietario tiene propiedad exclusiva de su lote. \\
Se requiere edificio ya construido para la afectación. & Puede estar construido o en proceso de construcción. \\
Típicamente: edificios en altura, plantas divididas. & Típicamente: lotes individuales con cerramiento perimetral. \\
La unidad privativa es un ``cubo de aire'' (espacio aéreo delimitado). & La unidad privativa incluye la tierra (el lote físico). \\
Destino predominantemente residencial o comercial en altura. & Cualquier destino: residencial, industrial, náutico, deportivo. \\
Regulación: arts.~2037--2072 CCyC. & Regulación: arts.~2073--2086 CCyC (propiedad horizontal especial). \\
\bottomrule
\end{tabularx}
\end{table}

\section{Partes comunes y privativas}

Son \textbf{necesariamente comunes} las partes del terreno destinadas a vías de circulación, acceso común y comunicación entre los distintos sectores; las destinadas a actividades deportivas, recreativas y sociales; y todo lo que el reglamento determine como común.

Son \textbf{partes privativas} las unidades funcionales, que pueden estar construidas o en proceso de construcción, y que deben tener independencia funcional según su destino y salida a la vía pública. Dada la naturaleza del emprendimiento, esa salida suele ser indirecta, a través del cerramiento perimetral.

\begin{quote}
\textbf{Artículo 2079 CCyC (cosas y partes necesariamente comunes).} \textit{Son necesariamente comunes o de uso común las partes y lugares del terreno destinadas a vías de circulación, acceso y comunicación, áreas específicas destinadas al desarrollo de actividades deportivas, recreativas y sociales, instalaciones y servicios comunes, y todo otro bien afectado al uso comunitario, calificado como tal por el reglamento de propiedad y administración que regula el emprendimiento.}
\end{quote}

\section{Persona jurídica, reglamento y normativa aplicable}

El conjunto inmobiliario conforma una \textbf{persona jurídica}, como ocurría con el consorcio de propietarios, y se constituye con un \textbf{reglamento} que define:

\begin{itemize}
    \item cuáles son las partes propias y cuáles las comunes;
    \item el funcionamiento y el régimen de asambleas;
    \item el régimen disciplinario;
    \item los gastos y la forma de calcularlos;
    \item el porcentual que corresponde a cada propiedad para contribuir a los gastos comunes.
\end{itemize}

Además de las normas del Código, se aplican todas las \textbf{normas urbanísticas} de cada jurisdicción. Son normas administrativas: no es posible construir un nuevo emprendimiento y someterlo al régimen de conjuntos inmobiliarios en cualquier lugar y en cualquier tierra. La localización, los límites perimetrales y los aspectos edilicios son regulados íntegramente por las normas provinciales y municipales, según se indica en el artículo 2079.
% TODO: contenido incompleto o ambiguo en las notas originales (las notas remiten al art. 2079 para las normas urbanísticas, mientras que el texto transcripto de ese artículo trata de las partes necesariamente comunes)

\section{Facultades y obligaciones del propietario}

El propietario debe:

\begin{itemize}
    \item ejercer su derecho conforme a la normativa, al reglamento y a las normas de convivencia;
    \item respetar los valores en materia de construcción, paisajísticos, arquitectónicos y ecológicos;
    \item pagar las \textbf{expensas} en la proporción que establece el reglamento.
\end{itemize}

El reglamento puede determinar además otras contribuciones por el uso de servicios; hay conjuntos donde se practica polo u otras actividades que implican costos adicionales.

\subsection{Restricciones reglamentarias propias de los conjuntos inmobiliarios}

En los conjuntos inmobiliarios existen restricciones que normalmente no se observan en la propiedad horizontal clásica:

\begin{itemize}
    \item restricciones sobre el régimen de invitados;
    \item reglamentación sobre quiénes pueden usar la pileta, las canchas de tenis y otros servicios comunes;
    \item posibilidad de cobrar un canon adicional por el uso de ciertas instalaciones;
    \item restricciones arquitectónicas, paisajísticas y ecológicas en las construcciones.
\end{itemize}

\subsection{Límite del reglamento: la transmisión de las unidades}

\begin{important}
El reglamento puede \textbf{restringir} el uso de los servicios y espacios comunes, pero \textbf{nunca puede prohibir} que un propietario venda o transmita su unidad, porque sería limitar la transmisión inmobiliaria. Sí puede establecer un \textbf{derecho de preferencia} a favor del consorcio o del resto de los propietarios (es decir, que puedan igualar la oferta antes que un tercero).
\end{important}

% ----------------------------------------------------------
\chapter[Conclusiones]{Conclusiones}

Los puntos centrales de este tema son los siguientes:

\begin{itemize}
    \item \textbf{Todo conjunto inmobiliario constituido a partir del 31 de agosto de 2015 solo puede organizarse como derecho real.}
    \item Los conjuntos preexistentes organizados como derechos personales o mixtos deben adecuarse, aunque el Código no fijó plazo ni sanción.
    \item La adecuación es costosa y compleja, por lo que durante muchos años coexistirán ambos tipos de organización.
    \item El profesional del derecho deberá identificar bajo qué régimen opera cada emprendimiento para asesorar correctamente.
\end{itemize}

En el ejercicio profesional se advertirá con facilidad que coexisten distintas formas de organización, tanto en la esfera de los derechos personales como en la de los derechos reales.

% ==========================================================
% CORNELL
% ==========================================================
\chapter[Cuadro Cornell]{Cuadro Cornell: Conjuntos inmobiliarios}

\renewcommand{\arraystretch}{1.25}
\begin{longtable}{
    >{\raggedright\arraybackslash}p{0.24\textwidth}
    >{\raggedright\arraybackslash}p{0.68\textwidth}
}
\toprule
\textbf{Preguntas / palabras clave} & \textbf{Notas / respuestas} \\
\midrule
\endfirsthead

\toprule
\textbf{Preguntas / palabras clave} & \textbf{Notas / respuestas} \\
\midrule
\endhead

\midrule
\multicolumn{2}{r}{\small\textit{Continúa en la página siguiente}} \\
\endfoot

\bottomrule
\endlastfoot

\textbf{¿Qué son los conjuntos inmobiliarios?} &
Son emprendimientos urbanísticos ---como countries, barrios cerrados, parques industriales o clubes náuticos--- que combinan partes privativas (el lote de cada propietario) con partes comunes de uso compartido, dentro de un cerramiento perimetral. El Código Civil y Comercial los regula como un derecho real de propiedad horizontal especial. \\[0.6em]

\textbf{¿Cuál era el problema de organizarlos como derechos personales?} &
Cuando la tierra pertenecía a una sociedad anónima y los compradores solo adquirían acciones, cualquier ejecución judicial o quiebra de esa empresa implicaba la subasta del terreno junto con todo lo construido. El adquirente perdía su inmueble sin protección jurídica real. \\[0.6em]

\textbf{¿Por qué era inconstitucional la norma bonaerense?} &
Las provincias no pueden legislar en materia de propiedad (es competencia del Congreso Nacional). Además, la propiedad horizontal exige que el terreno sea común a todos los copropietarios, lo cual es incompatible con la naturaleza de un country, donde cada propietario quiere su lote exclusivo. \\[0.6em]

\textbf{Artículo 2075 CCyC: ¿qué establece?} &
Dispone que todos los conjuntos inmobiliarios constituidos a partir del 31 de agosto de 2015 deben organizarse únicamente como derechos reales. Además, ordena que los emprendimientos preexistentes organizados como derechos personales o mixtos se adecuen al nuevo régimen, aunque el Código no fijó plazo ni sanción para esa adecuación. \\[0.6em]

\textbf{¿Qué diferencia al conjunto inmobiliario de la propiedad horizontal clásica?} &
En la propiedad horizontal tradicional, el terreno es siempre un bien común de todos los copropietarios. En el conjunto inmobiliario, cada propietario tiene la propiedad exclusiva de su lote (la tierra que le corresponde). La parte común son los espacios de circulación, recreación e infraestructura compartida. \\[0.6em]

\textbf{¿Cuáles son las partes necesariamente comunes?} &
Las vías de circulación y acceso común, los espacios de comunicación entre sectores y las instalaciones destinadas a actividades deportivas, recreativas y sociales. También son comunes las áreas que el reglamento designe expresamente para uso compartido de todos los propietarios. \\[0.6em]

\textbf{¿Qué ventajas otorga el derecho real al adquirente?} &
Otorga oponibilidad erga omnes (frente a todos), permite afectar el inmueble a la protección de la vivienda familiar, reconoce el derecho de habitación vitalicio y gratuito al cónyuge o conviviente supérstite y habilita la constitución de hipoteca y otros derechos reales sobre el bien. \\[0.6em]

\textbf{¿Qué puede y qué no puede restringir el reglamento?} &
Puede establecer restricciones sobre el régimen de invitados, el uso de instalaciones comunes (pileta, canchas), el tipo de construcciones y los valores arquitectónicos y paisajísticos. También puede reconocer un derecho de preferencia a favor del consorcio. No puede prohibir la transmisión de las unidades. \\[0.6em]

\textbf{¿Por qué es compleja la adecuación de los emprendimientos anteriores?} &
Porque requiere mensura y planos aprobados, escritura pública e inscripción registral para cada unidad. Estos trámites generan costos muy elevados. Por esa razón, el Código no impuso un plazo ni una sanción: muchos emprendimientos anteriores a 2015 seguirán funcionando como derechos personales durante años. \\[0.6em]

\textbf{¿Para qué sirve este conocimiento en la vida profesional?} &
El abogado o escribano debe saber identificar bajo qué régimen está organizado cada emprendimiento para asesorar correctamente: qué derechos tiene el cliente, cómo proteger el inmueble, si puede hipotecarlo, afectarlo a bien de familia o qué ocurre en caso de quiebra del desarrollista. \\

\midrule

\multicolumn{2}{p{\dimexpr0.92\textwidth+2\tabcolsep\relax}}{%
\textbf{Resumen:} Los conjuntos inmobiliarios representan una de las innovaciones más importantes del Código Civil y Comercial en materia de derechos reales. Antes de su sanción, emprendimientos como countries y barrios cerrados funcionaban en la esfera de los derechos personales, con graves riesgos para los adquirentes: la tierra pertenecía a una persona jurídica cuya quiebra podía provocar la pérdida total de lo construido. La doctrina impulsó su incorporación como derecho real, y el artículo 1887 del CCyC los reconoce expresamente. Como propiedad horizontal especial, cada propietario tiene el dominio exclusivo de su lote, mientras los espacios comunes permanecen en indivisión forzosa. El régimen se completa con un reglamento que organiza la convivencia, las expensas y las restricciones permitidas. La obligación de adecuación de los emprendimientos preexistentes plantea un desafío práctico: sin plazo ni sanción establecidos, la coexistencia de ambos regímenes será una realidad durante muchos años, y el profesional del derecho deberá saber identificarlos y asesorar en consecuencia.
} \\
\end{longtable}

\end{document}
